\documentclass[aspectratio=169, 10pt]{beamer}
\usepackage{beamer_theme_cienciadados2026}

\title[Aula 11 - NumPy Vetorial]{Ciência dos Dados I: Análise Exploratória}
\subtitle{Aula 11: Computação Científica de Alta Performance com NumPy}
\author{Prof. Dr. Ney Lemke}
\institute{UNESP -- Instituto de Biociências de Botucatu}
\date{Versão 2026}

\begin{document}

\frame{\titlepage}

\begin{frame}{Sumário da Aula}
  \tableofcontents
\end{frame}

\section{Por que NumPy?}

\begin{frame}{Estrutura Interna do \texttt{ndarray}}
  Listas padrão de Python guardam ponteiros para objetos genéricos espalhados na memória heap.
  \begin{block}{O Array NumPy (\texttt{np.ndarray})}
    \begin{itemize}
      \item Bloco contíguo de bytes em memória RAM (escrito em C/Fortran).
      \item Tipo homogêneo e fixo (\texttt{dtype}: \texttt{float64}, \texttt{int32}, etc.).
      \item Permite ao processador utilizar instruções SIMD (Single Instruction, Multiple Data) e vetorização no nível do hardware.
      \item Acelerações típicas de \textbf{10x a 100x} em relação a loops \texttt{for} nativos de Python!
    \end{itemize}
  \end{block}
\end{frame}

\section{Vetorização e Universal Functions}

\begin{frame}[fragile]{Vetorização: Eliminando Laços Explícitos}
  Em vez de iterar elemento por elemento, aplicamos a operação ao vetor inteiro:
  \begin{lstlisting}
import numpy as np

arr = np.linspace(0, 10, 1_000_000)

# Vetorizacao com Universal Functions (ufuncs)
seno = np.sin(arr)
potencia = arr ** 2
  \end{lstlisting}
  \begin{exampleblock}{Código Limpo e Rápido}
    A sintaxe aproxima o código Python diretamente da notação matemática da álgebra linear.
  \end{exampleblock}
\end{frame}

\section{Regras de Broadcasting}

\begin{frame}[fragile]{Operando em Arrays com Formatos Distintos}
  Broadcasting descreve como o NumPy trata arrays com shapes diferentes durante operações aritméticas:
  \begin{itemize}
    \item Se os arrays não possuem o mesmo número de dimensões, adiciona-se 1 à esquerda.
    \item Arrays com tamanho 1 em uma dimensão são "esticados" para coincidir com a dimensão maior sem alocar memória adicional.
  \end{itemize}
  \begin{lstlisting}
matriz = np.ones((3, 3)) # Shape (3, 3)
linha = np.array([1, 2, 3]) # Shape (3,) -> tratado como (1, 3)
resultado = matriz + linha # Shape final (3, 3)
  \end{lstlisting}
\end{frame}

\begin{frame}{Resumo da Aula}
  \begin{itemize}
    \item NumPy é o alicerce numérico de todo o ecossistema científico (Pandas, SciPy, Scikit-Learn, PyTorch).
    \item Vetorização e broadcasting viabilizam simulações em larga escala.
    \item \textbf{Prática}: Criação de arrays, benchmarking com \texttt{\%timeit} e normalização de dados.
  \end{itemize}
\end{frame}

\end{document}
